
Then, by IV (2.12.4), one has:
@@M0@@
the last equality following from the fact that for @@M1@@ one has @@M2@@ (VI 2.8). Suppose that @@M3@@. Then the preceding argument shows that @@M4@@; applying the same reasoning to the elements @@M5@@ for @@M6@@, one deduces by induction that the @@M7@@ are all zero, whence injectivity.

\sgaStableTarget{sga6.e09.para.3-3}3.3. Let @@M8@@ still be a noetherian scheme, and @@M9@@ a locally free @@M10@@-Module of rank @@M11@@. Denote by @@M12@@ the vector bundle (contravariant in @@M13@@!) associated with @@M14@@ (SGA 3 I 4.6), regarded as a group scheme over @@M15@@. Let @@M16@@ be a torsor under @@M17@@; we wish to compute @@M18@@.

Recall how one can embed @@M19@@ in a projective bundle over @@M20@@. Torsors under @@M21@@ correspond to extensions @@M22@@ of @@M23@@ by @@M24@@. Indeed, given an extension:
\sgaStableTarget{sga6.e09.eq.3-3-1}@@M25@@
one obtains a torsor under @@M26@@ by considering the homomorphism @@M27@@ defined by @@M28@@, the section of @@M29@@ defined by @@M30@@, and taking @@M31@@, the actions of @@M32@@ on @@M33@@ being induced by those of @@M34@@ on @@M35@@.

Conversely, let @@M36@@ be a torsor under @@M37@@. Denote by @@M38@@ the sheaf of germs of homomorphisms from @@M39@@ to @@M40@@; it is a principal homogeneous sheaf under the sheaf of germs of homomorphisms from @@M41@@ to @@M42@@, which identifies with the dual @@M43@@ of @@M44@@. In the classical way this yields an extension of @@M45@@ by @@M46@@, whence, by passing to duals, the sequence (3.3.1).

